\section{HyperGraphRAG}

\begin{frame}{HyperGraphRAG: Metadata}
\small
\textbf{HyperGraphRAG: Retrieval-Augmented Generation via Hypergraph-Structured Knowledge Representation}

\paperinfo
{NeurIPS 2025, 39th Conference on Neural Information Processing Systems}
{OpenAI \texttt{text-embedding-3-small} for entities and hyperedges}
{OpenAI GPT-4o-mini for extraction and generation}
{Domain RAG with n-ary knowledge hypergraph retrieval}

\begin{block}{Key insights / contributions}
\scriptsize
\begin{itemize}
  \item First hypergraph-based RAG method for representing n-ary factual relations.
  \item Combines knowledge hypergraph construction, hypergraph retrieval, and hypergraph-guided generation.
  \item Improves domain QA beyond chunk-based and binary KG-based RAG baselines.
\end{itemize}
\end{block}
\end{frame}

\begin{frame}{HyperGraphRAG: Problem and Failure Mode}
\scriptsize
\begin{columns}[T,onlytextwidth]
  \begin{column}{0.48\textwidth}
    \begin{block}{Problem}
      Many domain questions need one fact that binds several entities at once, e.g., patient type, lab range, diagnosis, and treatment condition.
    \end{block}
    \begin{block}{Why standard RAG fails}
      Chunk retrieval scores passages independently. It may retrieve a nearby paragraph but cannot explicitly preserve which entities participate in the same clinical or legal fact.
    \end{block}
  \end{column}
  \begin{column}{0.48\textwidth}
    \begin{block}{Why binary GraphRAG fails}
      Binary edges fragment an n-ary statement into many pairwise triples, creating sparse and incomplete evidence during retrieval.
    \end{block}
    \begin{block}{Method response}
      HyperGraphRAG stores each multi-entity fact as a hyperedge, retrieves both entities and hyperedges, then fuses them with chunks for generation.
    \end{block}
  \end{column}
\end{columns}
\end{frame}

\pipelineframe{HyperGraphRAG}{images/hypergraph_rag_pipeline.png}{0.98}

\begin{frame}{HyperGraphRAG: Algorithm Deep Dive}
\scriptsize
\begin{enumerate}
  \item \textbf{Fragment and extract:} the LLM segments documents into knowledge fragments and extracts entities plus n-ary relational descriptions.
  \item \textbf{Store as bipartite hypergraph:} entity nodes connect to hyperedge fact nodes; entity, hyperedge, and chunk vectors are indexed separately.
  \item \textbf{Query parsing:} the LLM extracts key entities from the question; query/entity and query/hyperedge similarities produce initial candidates.
  \item \textbf{Hypergraph expansion:} retrieved entities pull in incident hyperedges; retrieved hyperedges expose their participant entities and source chunks.
  \item \textbf{Knowledge fusion:} final context mixes structured n-ary facts with chunk evidence, then GPT-4o-mini generates the answer.
\end{enumerate}
\begin{block}{Core intuition}
The retrieval unit is not a single passage or binary edge; it is a multi-entity fact that preserves the relation scope needed by domain questions.
\end{block}
\end{frame}

\begin{frame}{HyperGraphRAG: Concrete Example}
\scriptsize
\begin{columns}[T,onlytextwidth]
  \begin{column}{0.47\textwidth}
    \begin{block}{Question}
      Which treatment is suitable for a male hypertensive patient with mild serum creatinine elevation?
    \end{block}
    \begin{block}{Failure case}
      A chunk may mention only the diagnosis, while binary triples split one clinical fact into weak edges:
      \texttt{patient--has--hypertension}, \texttt{patient--has--male}, \texttt{patient--has--creatinine}.
    \end{block}
  \end{column}
  \begin{column}{0.49\textwidth}
    \begin{block}{Hypergraph view}
      One hyperedge stores the whole n-ary fact:
      \begin{align*}
      e_h = \{&\text{patient}, \text{male}, \text{hypertension},\\
              &\text{creatinine range}, \text{treatment note}\}
      \end{align*}
      The retriever can match any subset of entities and still recover the complete fact.
    \end{block}
  \end{column}
\end{columns}
\end{frame}

\begin{frame}{HyperGraphRAG: Main Results}
\scriptsize
\begin{block}{Evaluation setup}
Five domains: Medicine, Agriculture, CS, Legal, and Mix. Metrics: answer F1, retrieval similarity (R-S), and LLM-judged generation evaluation (G-E).
\end{block}
\centering
\resizebox{0.96\linewidth}{!}{
\begin{tabular}{lccc}
\toprule
Method & Avg. F1 & Avg. R-S & Avg. G-E \\
\midrule
StandardRAG & 32.96 & 51.14 & 58.85 \\
Best binary graph baseline & 26.27 & 45.24 & 53.03 \\
HyperGraphRAG & \textbf{38.61} & \textbf{64.25} & \textbf{61.52} \\
\bottomrule
\end{tabular}
}
\vspace{0.2cm}
\begin{block}{Takeaway}
The largest gain is retrieval quality: n-ary hyperedges recover more relevant gold knowledge than binary graph structures, which then improves answer accuracy and generation quality.
\end{block}
\end{frame}
